
\documentclass[12pt]{article} % use larger type; default would be 10pt
\usepackage[square]{natbib}
\usepackage[utf8]{inputenc} % set input encoding (not needed with XeLaTeX)
\usepackage{amsmath,amssymb}
\DeclareMathOperator{\E}{\mathbb{E}}
\usepackage{color}
\usepackage{pdfpages}
%%% PAGE DIMENSIONS
\usepackage{geometry} % to change the page dimensions
%\geometry{a4paper} % or letterpaper (US) or a5paper or....
\geometry{margin=1in} % for example, change the margins to 2 inches all round
% \geometry{landscape} % set up the page for landscape
%   read geometry.pdf for detailed page layout information
 \newcommand{\eps}{\varepsilon}
\usepackage{graphicx} % support the \includegraphics command and options

% \usepackage[parfill]{parskip} % Activate to begin paragraphs with an empty line rather than an indent

%%% PACKAGES
\usepackage{booktabs} % for much better looking tables
\usepackage{array} % for better arrays (eg matrices) in maths
\usepackage{paralist} % very flexible & customisable lists (eg. enumerate/itemize, etc.)
\usepackage{verbatim} % adds environment for commenting out blocks of text & for better verbatim
\usepackage{subfig} % make it possible to include more than one captioned figure/table in a single float
% These packages are all incorporated in the memoir class to one degree or another...

%%% HEADERS & FOOTERS
\usepackage{fancyhdr} % This should be set AFTER setting up the page geometry
\pagestyle{fancy} % options: empty , plain , fancy
\fancyhf{}
\renewcommand{\headrulewidth}{1pt} % customise the layout...
\lhead{STAT 429 Time Series Analysis}\chead{FALL 2020}\rhead{Instructor: Hyoeun Lee}
\lfoot{}\cfoot{}\rfoot{Page \thepage}




%%% SECTION TITLE APPEARANCE
\usepackage{sectsty}
\allsectionsfont{\sffamily\mdseries\upshape} % (See the fntguide.pdf for font help)
% (This matches ConTeXt defaults)

%%% ToC (table of contents) APPEARANCE
\usepackage[nottoc,notlof,notlot]{tocbibind} % Put the bibliography in the ToC
\usepackage[titles,subfigure]{tocloft} % Alter the style of the Table of Contents
\renewcommand{\cftsecfont}{\rmfamily\mdseries\upshape}
\renewcommand{\cftsecpagefont}{\rmfamily\mdseries\upshape} % No bold!


\definecolor{Red}{rgb}{1,0,0}
\newcommand{\Red}{\color{Red}}
\definecolor{Blue}{rgb}{0,0,1}
\newcommand{\Blue}{\color{Blue}}


%%%%%%%THEOREM 
\usepackage{thmtools,thm-restate}
\newtheorem{theorem}{Theorem}
\newtheorem{lemma}[theorem]{Lemma}
\newtheorem{proposition}[theorem]{Proposition}
\newtheorem{corollary}[theorem]{Corollary}
\newtheorem{definition}[theorem]{Definition}

\newenvironment{proof}[1][Proof]{\begin{trivlist}
\item[\hskip \labelsep {\bfseries #1}]}{\end{trivlist}}
%\newenvironment{definition}[1][Definition]{\begin{trivlist}
%\item[\hskip \labelsep {\bfseries #1}]}{\end{trivlist}}
\newenvironment{example}[1][Example]{\begin{trivlist}
\item[\hskip \labelsep {\bfseries #1}]}{\end{trivlist}}
\newenvironment{remark}[1][Remark]{\begin{trivlist}
\item[\hskip \labelsep {\bfseries #1}]}{\end{trivlist}}

\newcommand{\qed}{\nobreak \ifvmode \relax \else
      \ifdim\lastskip<1.5em \hskip-\lastskip
      \hskip1.5em plus0em minus0.5em \fi \nobreak
      \vrule height0.75em width0.5em depth0.25em\fi}
   
   
   
   
%   \usepackage{nopageno}   
 %%%%%%%%%%%%%%

%\linespread{2}





%%% END Article customizations



%%% The "real" document content comes below...


\date{}


%\title{Chapter 1}

\begin{document}




\textbf{\Large 1.2 Time Series Elements: Time Series Data}
%\textbf{\Large Chapter 1 Time Series Elements}\\
%\textbf{\large 1.1 Introduction}

\medskip


%\noindent
\fbox{%
    \parbox{0.9\textwidth}{%
 


\textbf{Learning outcome of Section 1.2}

\begin{itemize}
\item Understand/ Recall basic R commands for this course 
\item Plot Time Series Data with R
\item Analyze the Time Series Plot (visual inspection)
\item Understand concept of statistical model
\end{itemize}



    }%
}




\medskip
\medskip

\textbf{Recall}

\begin{itemize}
	\item Time series: a discrete stochastic process
	\item There is a probability law that governs every stochastic process
	\item Model: approximation of such a probability law.
\end{itemize}

\medskip
\textbf{Time Series Data}

\begin{itemize}
	\item Example 1.1, 1.3: JJ QEPS, DJIA\\ log-transformation, growth rate, volatility
	\item Example 1.2: Climate data\\ Description of trend, is the trend linear?
	\item Example 1.4: El Ni{\~n}o \\ Cycle and frequency
	\item Example 1.5: Predator-Prey Interactions \\  Relationship between two Time Series
\end{itemize}

\textbf{3 things important for modelling: Mean, Variance, Dependence}

many time series analysis comes to handling these 3 things well.

\hrulefill

Next: 1.3 Time Series models
\begin{itemize}
	\item Example of models for time series:
	\item white noise, smoothing, autoregression, random walk, signal+noise
\end{itemize}

\end{document}
